
No capítulo de Estado da Arte serão explorados e analisados trabalhos correlatos que se destacam na literatura científica em relação ao tema proposto. Este capítulo fornecerá uma base de revisões e críticas das abordagens existentes para a previsão utilizando aprendizado de máquina, com foco em incêndios florestais e relatando, especialmente, estudos que envolvem a análise empírica e estatística dos dados para correlacionar as ocorrências e os eventos que se fazem presente. Serão apresentadas contribuições significativas, metodologias inovadoras, e resultados obtidos por pesquisadores que exploraram o emprego de técnicas avançadas de inteligência artificial na antecipação destes eventos.

\section{Trabalhos Correlatos}

Os trabalhos de previsão de incêndios florestais dividem seus esforços em dois principais métodos de análise, que de um lado consistem na utilização de modelos matemáticos e estatíscos e do outro, no Aprendizado de Máquinas.

A previsão baseada em estatísticas ocorre por meio da correlação dos parâmetros metereológicos registrados nos incêndios históricos, tais como condições climáticas, terreno, temperatura média, pressão circunstancial e a frequência de incêndios que estes fatores causam ao ambiente. Pradeep et al. (citar) utilizaram um Sistema de Informação Geográfica (SIG) para delinear as zonas de risco de incêndios florestais no Parque Nacional Eravikulam, na Índia, por meio dos métodos de Processo Analítico Hierárquico (AHP) e a Razão de Frequência (FR), chegando a categorizar 72\% e 24\%, respectivamente, das incidências de incêndio na zona de alto risco dos mapas preparados. Gülçin et al. (citar) também elaboraram um mapa de risco de incêndio, dessa vez na província de Manisa, na Turquia, e utilizaram técnicas baseadas em Sensoriamento Remoto (SR) e SIG, reportando que 25,8\% da área de estudo apresentava risco alto ou muito alto de combustão espontânea.

A velocidade de propagação do fogo e a proporção do alcance do mesmo dependem fortemente da topografia e a arborização da região. Jin et al. (citar) realizaram um estudo sobre as características de incêndios florestais em duas províncias da Coréia do Norte (South Hamgyong e Gangwon) e concluíram que as árvores coníferas são muito mais suscetíveis à combustão em relação às demais plantas da região. Além disso, os resultados mostraram que a maioria dos grandes incêndios ocorreu ao longo de encostas abertas a favor do vento. Shi et al. (citar) realizaram um estudo para relacionar os danos ao corredor das linhas de energia de Hubei, na China, com as incidências de incêndios florestais nas suas proximidades. Os resultados foram obtidos a partir da implementação de um modelo de Regressão Logística para prever as ocorrências de incêndio e sua correlação com os desligamentos da linha, cruzando dados históricos e topográficos.

No que tange a previsão a partir de dados tabulares, os métodos de Aprendizado de Máquina que se mantém em evidência são os modelos mais tradicionais, como Regressão Logística, Floresta Aleatória e \textit{Suport Vector Machine} (SVM). Kalantar et al. (citar) previu a suscetibilidade de incêndios florestais na bacia hidrográfica de Chaloos Rood, no Irã, aplicando três modelos de aprendizado de máquina à um conjunto de 14 variáveis preditoras, que derivaram das características climáticas, ambientais e topográficas da região. Destaque para o modelo de \textit{boosted regression tree} (BRT), que superou os outros dois de forma geral, alcançando os maiores índices de Área Sob Curva (AUC) na métrica de desempenho. Bui et al. (citar) propôs a utilização de um modelo de aprendizado promissor para a classificação binária de suscetibilidade de fogo, chamado \textit{Kernel Logistic Regression} (KLR). A avaliação foi realizada nas proximidades do Parque Nacional Cat Ba, no Vietnã. De forma simplista, o algoritmo é capaz de mapear o conjunto de entrada de maneira linear e atribui um "peso" de influência para cada uma das características sobre o resultado final e, neste estudo, foi capaz de alcançar resultados extremamente satisfatórios. Para título de comparação, foi utilizada uma SVM de referência nesta mesma base dados e com os mesmos parâmetros do modelo proposto, e os resultados descriminaram uma leve vantagem no desempenho do modelo KLR.

Na província de Heilongjiang, Chang et al. (citar) utilizaram regressão logística para criar um mapa de probabilidade de ignição. Com uma proposta um pouco mais específica, este trabalho categorizou os incêndios entre antropogênicos (causados por humanos) e naturais. Dessa forma, a importância que as variáveis dependentes possuíam foi determinada a partir das correlações identificadas na incidência de cada um dos tipos de fogo. Concluiu-se na região estudada que os incêndios naturais possuem uma forte relação com a média anual de temperatura, enquanto que a variável mais influente para incêndios antropogênicos é o tipo de vegetação. Bisquert et al. (citar) classificaram os seus resultados em três níveis de suscetibilidade. Para isso, utilizou-se o modelo de regressão logística para identificar as variáveis de maior importância dentre os verdadeiros positivos e, posterior à essa classificação, estes dados foram inseridos em um algoritmo de Rede Neural Artificial, atribuindo três possíveis saídas aos rótulos do modelo (de menor à maior risco de fogo).

Naser et al. (citar) exploraram uma variedade de modelos de aprendizado de máquina, entre eles algoritmos de aprendizado profundo e árvore de decisão que apresentaram uma precisão superior a 80\% em bases de dados extremamente grandes. Além disso, os algoritmos genéticos também foram efetivos na previsão do tamanho esperado dos incêndios, também com margem de precisão acima de 80\%. A pesquisa concluiu que as performances atingidas tem uma ligação direta com a qualidade das observações dos incêndios que são documentados, tais como dados históricos, do clima e relevo. Malik et al. (citar) propuseram duas abordagens para o aprendizado de risco de incêndio. A primeira, a qual foi denominada modelo de conjunto, utiliza dois modelos base para segregar o conjunto de variáveis em escopos e, com os resultados alcançados nestes blocos, seguir para um novo treinamento. O modelo de conjunto se provou incapaz de derivar as relações entre as variáveis independentes para atingir resultados mais concisos. A segunda abordagem utilizou apenas um modelo que recebeu um conjunto completo dos dados selecionados, chamado de modelo combinado. Este segundo experimento revelou um agente capaz de deduzir efetivamente relações entre as variáveis independentes, uma vez que elas foram testadas em conjunto, e alcançou precisão superior a 90\%.

Pode-se segmentar o monitoramento de incêndios florestais em dois grandes espectros: a previsão, que determina as áreas de risco e antecede as medidas preventivas; e a detecção, que viabiliza o monitoramento em tempo real e visa distinguir os incêndios reais de maneira pontual. Pode-se dizer que as pesquisas voltadas para previsão e detecção estão fortemente ligadas enquanto finalidade (prevenção de incêndios), mas se dissociam nas metodologias, modelos e resultados esperados. Todos os estudos supracitados baseiam-se no reconhecimento estatístico de padrões históricos de queimadas, mas quando cita-se monitoramento de incêndios, não se pode deixar de lado estudos que visam reconhecer o fogo por meio da visão computacional. Fan et al. (citar) propuseram um modelo de Aprendizado Profundo através da combinação de duas arquiteturas tradicionais: o \textit{lightweigth vision transformer} (ViT) e as Redes Neurais Convolucionais (CNN's). O objetivo desta combinação era otimizar os resultados na detecção de incêndios florestais, diminuindo a complexidade computacional e conservando a perfomance dos algoritmos tradicionais. O SWVR, como foi chamado o modelo proposto, alcançou uma precisão de 79,6\% na identificação de incêndios reais. Além disso, foi capaz reduzir o custo computacional do modelo de referência em quase 37\%.

\subsection{Literatura recente (2024--2025)}

A última década consolidou o uso de \emph{ensembles} de árvores e métodos de interpretabilidade como prática essencial em estudos de suscetibilidade a incêndios florestais. Três tendências contemporâneas merecem destaque por nortearem as decisões metodológicas deste trabalho.

\paragraph{Interpretabilidade via SHAP e \emph{Permutation Importance}.} Cheerala \emph{et al.}~\cite{cheerala2025probabilistic} aplicaram \emph{Random Forest} combinado com SHAP~\cite{lundberg2017unified} em regiões de \emph{grassland} e \emph{forest} da Califórnia, obtendo AUC de \num{0,996} e \num{0,997} respectivamente, com identificação clara das \emph{features} dominantes (NDVI, EVI, histórico de fogo, umidade do combustível). Vasconcelos \emph{et al.}~\cite{forests2024cerradoamazon} identificaram o índice KBDI (\emph{Keetch-Byram Drought Index})~\cite{keetch1968drought} como o melhor preditor de fogo em Canaã dos Carajás (Pará), na zona de transição Cerrado--Amazônia. Quesada-Ruiz \emph{et al.}~\cite{npjhazards2025fire}, em \emph{npj Natural Hazards}, propuseram um modelo híbrido (componente dinâmica do ECMWF + \emph{Random Forest}) que utiliza o SPI~\cite{mckee1993spi} observado para prever anomalia de área queimada com até um mês de antecedência em \(\approx 68\%\) das áreas queimáveis brasileiras. Esses trabalhos motivam, na presente pesquisa, a adoção combinada de KBDI \emph{proxy}, VPD \emph{proxy}~\cite{seager2015climatology} e SPI multi-escala como \emph{features} Tier~1 do \emph{pipeline} (Capítulo~\ref{cap:metodologia}).

\paragraph{Validação temporal e \emph{drift} climático.} Bergmeir \& Benítez~\cite{bergmeir2012cv} alertam que o uso de \emph{cross-validation} estratificada padrão em séries temporais introduz viés otimista, sendo necessário utilizar \emph{rolling-origin} ou \emph{TimeSeriesSplit}. Esta recomendação é tomada como princípio metodológico no Capítulo~\ref{cap:resultados} deste trabalho. A relevância prática desse cuidado é amplificada por Silva-Junior \emph{et al.}~\cite{silvajunior2025amazon}, que documentam 2024 como o pior ano de degradação por fogo na Amazônia em mais de duas décadas (\SI{3,3}{\mega\hectare}), justificando o uso de janelas temporais recentes (2021--2023) como teste mais desafiador --- período em que os modelos devem provar robustez sob \emph{drift} climático real.

\paragraph{\emph{Pseudo-labeling} para imputação climática.} Lopez-Garcia \emph{et al.}~\cite{lopezgarcia2024spatiotemporal} aplicaram \emph{pseudo-labeling}~\cite{lee2013pseudo} para imputar dados SMAP em \emph{grids} tropicais, com regressor supervisionado atingindo \(R^2 > 0{,}90\) em \emph{holdout}. A mesma estratégia é adotada neste projeto (Seção~\ref{sec:metodologia_pseudo_label}) para imputar a variável \texttt{Umidade} nos \num{763489} registros não enriquecidos via NASA POWER, com cuidados explícitos quanto ao viés de confirmação relatados por Arazo \emph{et al.}~\cite{arazo2020pseudo}.

\paragraph{Fusão estação--grade, \emph{fire weather} e ambiguidade operacional.} Revisões recentes enfatizam variáveis dinâmicas de combustível e meteorologia de superfície (precipitação local, déficit de umidade, temperatura máxima, vento) como condicionantes de propagabilidade, frequentemente sintetizadas em índices de \emph{fire weather}. Em paralelo, a literatura de observação da Terra discute quantificação de incerteza preditiva --- ainda raramente operacionalizada em painéis acadêmicos. Estas tendências motivam extensões ao \emph{pipeline} de serviço documentadas na Seção~\ref{sec:res_estudo_caso_pium} (subsubseção~M4): busca hierárquica de estações INMET com rótulo de representatividade, vento WIS2 alinhado à precipitação da mesma estação, variáveis NASA POWER adicionais expostas na API, indicador escalar de \emph{fire weather} e métricas FNR/FPR na auditoria contra FIRMS --- implementadas de modo a não alterar o pré-processador de treino até um retreino controlado, evitando \emph{train/serve skew}.

Em resumo, a bibliografia discutida destaca o amplo espectro de abordagens e tecnologias disponíveis para mitigar o impacto devastador dos incêndios florestais. Desde métodos tradicionais, como modelos baseados em dados históricos e estatística, até avanços recentes em inteligência artificial e aprendizado de máquina, a pesquisa nesse campo está constantemente evoluindo. Destaca-se, nesta área, a importância de uma base de dados sólida e composta por múltiplas fontes, incluindo informações meteorológicas, dados geoespaciais e situação climática, que podem contribuir fundamentalmente para melhorar a precisão das previsões. Além disso, apesar de não ser parte do escopo desta pesquisa, técnicas avançadas, como o uso de redes neurais convolucionais e algoritmos de aprendizado profundo, tem entregue grandes expectivas na identificação precoce de áreas queimadas e na mensuração de propagação do fogo.

